%%%%%%%%%%%%%%%%%%%%%%%%%%%%%%%%%%%%%%%%%%%%%%%%%%%%%%%%%%%%%%%%%%%%%%
% Lớp 11 | Đạo hàm | Đạo hàm cấp hai
% Dạng 1D7?3-3: Toán thực tế và Ý nghĩa Vật lý của đạo hàm cấp hai
% Số câu: 23  (sinh tự động bởi tach_cau.py; từ nay thầy sửa trực tiếp file này)
%%%%%%%%%%%%%%%%%%%%%%%%%%%%%%%%%%%%%%%%%%%%%%%%%%%%%%%%%%%%%%%%%%%%%%

%%%%%%%%%%%%%%%% TRẮC NGHIỆM 4 PHƯƠNG ÁN %%%%%%%%%%%%%%%%

% nguồn: Ok_Toan11_TracNghiem_Lan2
\begin{ex}%[1D7H3-3]%[ma:6795f112]
	Phương trình  chuyển động của một chất điểm được biểu thị bởi công thức $$s(t)=4-2t+4t^2+2t^3,$$ trong đó $t > 0$ và $t$ tính bằng giây (s), $s(t)$ tính bằng mét (m). Tìm gia tốc $a$ của chất điểm tại thời điểm $t=5$ (s).
	 \choice
		{\True  $a=68$ m/s$^2$}
		{$a=115$  m/s$^2$}
		{$a=100$  m/s$^2$}
		{$a=225$ m/s$^2$} 
	 \loigiai{
		Theo ứng dụng đạo hàm của hàm số ta có $v(t)=s'(t)=-2+8t+6t^2$ và  $a(t)=v'(t)=8+12t$.\\
		Vậy gia tốc của chất điểm tại thời điểm $t=5$ (s)  là $a(5)=68$ m/s$^2$.
	} 
\end{ex}


%%%%%%%%%%%%%%%% ĐÚNG/SAI %%%%%%%%%%%%%%%%

% nguồn: Ok_dungsai_lop11_Lan1
\begin{ex}%[1D7H3-3]%[ma:e370899e]
 Một chất điểm chuyển động có phương trình $S(t) = t^3 - 3t^2 - 9t + 2$; trong đó $S$ tính bằng mét (m), $t$ tính bằng  giây (s). Xét tính Đúng - Sai của các mệnh đề sau:
 \choiceTF 
  {\True Đạo hàm của $S(t)$ tại thời điểm $t=t_0$ chính là vận tốc tức thời của chất điểm tại thời điểm $t=t_0$} 
  {Vận tốc tức thời của chất điểm đó tại thời điểm $t=5$ (giây)  là $27$ m/s}
  {\True Gia tốc của chất điểm đó tại thời điểm vận tốc bằng  $0$ là $12$ m/s$^2$}
  {Thời điểm khi $t=6$ (giây) thì vận  tốc của chất điểm đạt giá trị nhỏ nhất.}
  \loigiai{
  \begin{itemchoice}
   \itemch Đúng. 
   \itemch  Sai.\\
   Ta có $v(t)=S'(t) = 3t^2 - 6t -  9$. \\
   Vận  tốc tức thời của chất điểm đó tại thời điểm $t=5$ (giây) là $v(t)=36$ m/s.
   \itemch Đúng.\\ 
   Ta có $S'(t)  = 3t^2 - 6t - 9 = 0 \Leftrightarrow \hoac{&t = 3 \, (\text{nhận}) \\&t = -1 \, (\text{loại}).}$\\
   $a(t)=S''(t) = 6t - 6$.\\ 
   Gia tốc của  chất điểm tại thời điểm vận tốc bằng $0$ là $a(3)=S''(3) = 12$ m/s$^2$.
    \itemch Sai.\\
   Ta có $v(t)=S'(t) = 3t^2 - 6t - 9 = 3(t-1)^2+6 \ge 6$.  Vậy vận tốc nhỏ nhất là $6$ m/s khi t=1 (giây).
  \end{itemchoice}
 } 
\end{ex}

% nguồn: Ok_dungsai_lop11_Lan2
\begin{ex}%[1D7H3-3]%[ma:6f633379]
	Chuyển động của một vật có phương trình $s(t)=4\cos\left(2\pi t-\dfrac{\pi}{12}\right)$ m, với $t$  là thời gian tính bằng giây.
	 \choiceTF
		{\True $s'(t)=-8\pi\sin\left(2\pi t-\dfrac{\pi}{12}\right)$} 
		 {$s''(t)=16\pi^2\cos\left(2\pi t-\dfrac{\pi}{12}\right)$}
		{\True Vận tốc  của vật tại thời điểm khi $t=5$ (s) xấp xỉ $6{,}505$ m/s}
		{Gia tốc của vật tại  thời điểm khi $t=5$ (s) xấp xỉ $152{,}533$ m/s$^2$}
	\loigiai{ 
		\begin{itemchoice}
			\itemch Ta  có $s'(t)=-8\pi\sin\left(2\pi t-\dfrac{\pi}{12}\right)$.
			\itemch Ta có $s''(t)=-16\pi^2\cos\left(2\pi t-\dfrac{\pi}{12}\right)$. 
			\itemch Vận tốc tức thời của vật tại thời  điểm $t=5$ (s) là
			$$s'(5)=-8\pi\sin\left(10\pi  -\dfrac{\pi}{12}\right)\approx 6{,}505\ \text{m/s}.$$
			\itemch Gia tốc tức thời của vật tại  thời điểm $t=5$ (s) là
			 $$s''(5)=-16\pi^2\cos\left(10\pi -\dfrac{\pi}{12}\right)\approx-152{,}533 \ \mathrm{m/s}^2.$$
		\end{itemchoice}
	 }
\end{ex}

% nguồn: Ok_dungsai_lop11_Lan1
\begin{ex}%[1D7V3-3]%[ma:05eca63a]
  Một vật chuyển động trên đường thẳng được xác định bởi công thức $s(t)=t^3-3 t^2+7 t-2$, trong đó $t>0$ và tính bằng giây và $s$ là quãng đường chuyển động được của vật trong $t$ giây tính bằng mét. Khi đó
  \choiceTF
  {\True Tốc  độ của vật tại thời điểm $t=2$ là $v(2)=s^{\prime}(2)=7$ m/s}
  {\True Gia tốc của vật tại thời điểm $t=2$ là $a(2)=v^{\prime}(2)=6$  m/s$^2$}
  {Gia tốc của vật tại thời điểm mà vận tốc  của chuyển động bằng $16$ m/s$^2$ là $ 10$ m/s$^2$}
  {\True Thời điểm $t=1$ giây tại đó vận tốc của chuyển động đạt giá  trị nhỏ nhất}
  \loigiai{
  Ta có $s^{\prime}(t)=3 t^2-6 t+7$ và  $s^{\prime\prime}(t)=6 t-6$.
  \begin{itemchoice}
   \itemch Vận tốc của vật  tại thời điểm $t=2$ là $v(2)=s^{\prime}(2)=3\cdot 2^2-6\cdot 2+7=7$ (m/s).
   \itemch Gia tốc của vật tại thời điểm  $t=2$ là $a(2)=v^{\prime}(2)=s^{\prime\prime}(2)=6\cdot 2-6=6$ (m/s$^2$).
   \itemch Vận tốc của chuyển động bằng $16$  m/s$^2$ tại thời điểm $t$ nghĩa là\\
    $ v(t)=s^{\prime}(t)=16 \Leftrightarrow 3t^2-6t+7=16\Leftrightarrow \hoac{& t=3\text{ (thoả mãn)} \\
   &  t=-1\text{(loại).}}$\\
   Gia tốc của vật tại thời điểm $t=3$ là $a(3)=v^{\prime}(3)=s^{\prime\prime}(3)=6\cdot 3-6=12$  (m/s$^2$).
   \itemch Vận tốc của chuyển động có phương trình $v(t)=3 t^2-6 t+7$ là một parabol, có đỉnh $S\left(-\dfrac{b}{2 a};-\dfrac{\Delta}{4 a}\right)\Rightarrow S(1; 4)$ và hệ số  $a=3>0$ nên hàm số có giá trị nhỏ nhất bằng $4$ tại $t=1$. Vậy tại thời điểm $t=1$ thì vận tốc của chuyển động đạt giá trị nhỏ nhất bằng $4$ m/s.
  \end{itemchoice}
 } 
\end{ex}

% nguồn: Ok_dungsai_lop11_Lan1
\begin{ex}%[1D7V3-3]%[ma:e52cb67e]
 Chuyển động của một vật có phương  trình $s(t)=4\cos\left(2\pi t-\dfrac{\pi}{12}\right)$ m, với $t$ là thời gian tính bằng giây. Khi đó
  \choiceTF
  {\True $s^{\prime}(t)=-8\pi\sin\left(2\pi t-\dfrac{\pi}{12}\right)$} 
  {$s^{\prime\prime}(t)=16\pi^2\cos\left(2\pi t-\dfrac{\pi}{12}\right)$} 
  {\True Vận tốc  của vật tại thời điểm khi $t=5$ s là
  $s'(5)=  \text{ m/s}$}
  {Gia tốc  của vật tại thời điểm khi $t=5$ s là $s''(5)\approx 152,533$ m/$s^2$}
  \loigiai{
  \begin{itemchoice}
    \itemch Đúng. Vì ta có $s^{\prime}(t)=-8\pi\sin\left(2\pi t-\dfrac{\pi}{12}\right)$ và $s^{\prime\prime}(t)=-16\pi^2\cos\left(2\pi t-\dfrac{\pi}{12}\right)$.
   \itemch Sai. Vì ta có $s^{\prime}(t)=-8\pi\sin\left(2\pi  t-\dfrac{\pi}{12}\right)$ và $s^{\prime\prime}(t)=-16\pi^2\cos\left(2\pi t-\dfrac{\pi}{12}\right)$.
   \itemch Đúng. Vì vận tốc tức thời của vật tại thời điểm $t=5$ s là  \\
    $s^{\prime}(5)=-8\pi\sin\left(10\pi-\dfrac{\pi}{12}\right)\approx 6,505$ (m/s).
   \itemch Đúng. Vì gia tốc tức thời của vật tại thời  điểm $t=5$ s là\\
    $s^{\prime\prime}(5)=-16\pi^2\cos\left(10\pi-\dfrac{\pi}{12}\right)\approx-152{,}533$ (m/s$^2$).
  \end{itemchoice}
 } 
\end{ex}

% nguồn: Ok_dungsai_lop11_Lan2
\begin{ex}%[1D7V3-3]%[ma:fc33e181]
	Phương trình chuyển động của một hạt dao động điều hòa là $s(t)=3\sin (2\pi t)$, trong đó $s$ tính bằng cm, $t$ là thời gian tính bằng giây. 
	 \choiceTF
		{Vận tốc tức thời của hạt tại thời  điểm $t=\dfrac{1}{6}$ (s) là $6\pi$ cm/s}
		{\True Gia tốc  tức thời của hạt tại thời điểm $t=\dfrac{1}{12}$ (s) là $-6\pi^2$ cm/s$^2$}
		{\True Trong khoảng thời gian từ $0{,}5$  giây đến $6{,}5$ giây, hạt đạt vận tốc cực đại $6$ lần}
		{\True Gia tốc  của hạt triệt tiêu khi nó đi qua điểm gốc}
	\loigiai{Vận tốc tức thời của chuyển động $v(t)=s'(t)=2\pi\cdot 3\cos (2\pi t)=6\pi\cos  (2\pi t)$.\\
		Gia tốc  tức thời của chuyển động $a(t)=v'(t)=-12\pi^2\sin (2\pi t)$.
		\begin{itemchoice}
			\itemch Vận tốc tức thời  của hạt tại thời điểm $t=\dfrac{1}{6}$ là $v\left(\dfrac{1}{6}\right)=6\pi\cdot\cos\left(\dfrac{\pi}{3}\right)=3\pi $.
			 \itemch Gia tốc tức thời của hạt tại thời điểm $t=\dfrac{1}{12}$ là $a\left(\dfrac{1}{12}\right)=-12\pi^2\sin\left(\dfrac{\pi}{6}\right)=-6\pi^2$.
			\itemch Ta có $v(t)=6\pi\cos (2\pi t)\leq 6\pi$,  $\forall t$.\\
			Đẳng thức xảy  ra khi $\cos (2\pi t)=1\Leftrightarrow 2\pi t=k2\pi , k\in\mathbb{Z}\Leftrightarrow t=k$.\\
			Vì $k\in\mathbb{Z}$ và $0{,}5\leq  t\leq 6{,}5$ nên $t\in\{1;2;3;4;5;6\}$.\\
			Vậy trong khoảng thời gian từ $0{,}5$ giây đến $6{,}5$ giây, hạt đạt vận tốc cực đại  $6$ lần.
			\itemch  Hạt đi qua điểm gốc khi $s(t)=0\Leftrightarrow 3\sin (2\pi t)=0\Leftrightarrow\sin (2\pi t)=0\Leftrightarrow a(t)=0$.
			Vậy gia tốc của hạt triệt tiêu khi nó đi  qua điểm gốc.
		\end{itemchoice}
	} 
\end{ex}

% nguồn: Ok_dungsai_lop11_Lan2
\begin{ex}%[1D7V3-3]%[ma:50b70786]
	Vị trí của một vật chuyển động thẳng được cho bởi phương trình $s(t)=\dfrac{1}{4}t^4-t^3+2t^2+2t$,  trong đó $s$ tính bằng mét (m) và thời gian $t$ tính bằng giây (s).
	\choiceTF 
		{\True Gia tốc của  vật tại thời điểm $t=2$ (s) là $4$ m/s$^2$}
		{Gia tốc  của vật tại thời điểm mà vận tốc đạt $34$ m/s là $24$ m/s$^2$}
		{Vận tốc  của vật tại thời điểm gia tốc bằng $13$ m/s$^2$ là $12$ m/s}
		{Vận tốc của vật  tại thời điểm gia tốc đạt giá trị nhỏ nhất trong $5$ giây đầu là $1$ m/s}
	 \loigiai{
		Vận tốc tức thời của chuyển động  $v(t)=s'(t)=t^3-3t^2+4t+2$.\\
		Gia tốc tức thời của chuyển động  $a(t)=v'(t)=3t^2-6t+4$.
		\begin{itemchoice}
			\itemch Gia tốc  của chuyển động tại thời điểm $t=2$ (s) là $a(2)=3\cdot 2^2-6\cdot 2+4=4$ m/s$^2$.
			\itemch Tại thời điểm vận tốc bằng $34$ m/s tức là $$v(t)=34\Leftrightarrow  t^3-3t^2+4t+2=34\Leftrightarrow t=4.$$
			Khi đó gia tốc của  vật là $a(4)=3\cdot 4^2-6\cdot 4+4=28$ m/s$^2$.
			\itemch Tại thời  điểm gia tốc bằng $13$ m/s$^2$ tức là $$a(t)=3t^2-6t+4=13\Leftrightarrow 3t^2-6t-9=0\Leftrightarrow\hoac{&t=3&&\\&t=-1&& (\text{loại}).}$$
			Khi đó vận tốc của vật là $v(3)=3^3-3\cdot 3^2+4\cdot  3+2=14$ m/s.
			\itemch Gia tốc tức thời của vật tại thời điểm $t$ là $a(t)=3t^2-6t+4=3(t-1)^2+1\geq 1,\ \forall  t\in[0;5]$.\\
			Do đó, trong $5$ giây đầu, gia tốc  đạt giá trị nhỏ nhất khi $t=1$.\\
			Khi đó vận tốc của vật là $v(1)=1^3-3\cdot 1^2+4\cdot 1+2=4$  m/s.
	\end{itemchoice}}
\end{ex}

% nguồn: Ok_dungsai_lop11_Lan2
\begin{ex}%[1D7V3-3]%[ma:b2b7c9fa]
	Phương trình chuyển động của một chất điểm được biểu thị bởi công thức $$s(t)=4-2t+4t^2-2t^3,$$ trong đó $t > 0$  và $t$ tính bằng giây, $s(t)$ tính bằng mét.
	 \choiceTF
		{Phương trình gia tốc của chất điểm là  $a(t)=12t+8$ m/s$^2$}
		{Gia tốc của chất điểm  tại thời điểm $t=5$ (s) là $52$ m/s$^2$}
		{\True Gia tốc của chất điểm triệt tiêu tại thời điểm vận tốc  lớn nhất}
		{\True Đến thời điểm $t=3$ (s) chất điểm chuyển động được quãng  đường xấp xỉ $24{,}59$ m}
	\loigiai{ 
		\begin{itemchoice}
			\itemch Ta có $v(t)=s'(t)=-6t^2+8t-2$  và $a(t)=v'(t)=-12t+8$.
			 \itemch Ta có gia tốc của chất điểm tại thời điểm $t=5$ (s) là $a(5)=-52$ m/s$^2$.
			 \itemch Vận tốc $v(t)$ của chất điểm là hàm số bậc hai có hệ số $a=-6<0$ nên chất điểm đạt vận tốc lớn nhất tại thời điểm $t=-\dfrac{b}{2a}=\dfrac{2}{3}$ (s).\\
			Khi đó gia tốc của chất điểm  là $a\left(\dfrac{2}{3}\right)=0$.
			\itemch  Ta có
			$v=0\Leftrightarrow-6t^2+8t-2=0\Leftrightarrow\hoac{&t=1\\ 
			 &t=\dfrac{1}{3}.}$\\
			Quãng đường đi được đến thời điểm $t=3$  (s) là
			$$\left|s\left(\dfrac{1}{3}\right)-s(0)\right|+\left|s(1)-s\left(\dfrac{1}{3}\right)\right|+|s(3)-s(1)|=\dfrac{8}{27}+\dfrac{8}{27}+24\approx 24{,}59\ \mathrm{m}.$$ 
		\end{itemchoice}
	 }
\end{ex}

% nguồn: Ok_dungsai_lop11_Lan2
\begin{ex}%[1D7V3-3]%[ma:3dee63c9]
	Một vật chuyển động trên đường thẳng được xác định bởi công thức $s(t)=t^3-3t^2+7t-2$, trong đó $t>0$ và tính bằng giây và $s$ là quãng đường chuyển động được của vật trong $t$ giây tính  bằng mét.
	 \choiceTF
		{\True Tốc độ của vật tại thời điểm $t=2$  là $7$ m/s}
		{\True Gia tốc của vật tại thời điểm $t=2$  là $6$ m/s$^2$}
		{Gia tốc của  vật tại thời điểm mà vận tốc của chuyển động bằng $16$ m/s$^2$ là $10$ m/s$^2$}
		{\True Thời điểm $t=1$ (giây) tại đó vận tốc của chuyển động đạt  giá trị nhỏ nhất}
	\loigiai{ 
		Ta có $s'(t)=3t^2-6t+7$ và  $s''(t)=6t-6$.
		\begin{itemchoice}
			\itemch Vận tốc của vật  tại thời điểm $t=2$ là $v(2)=s'(2)=3\cdot 2^2-6\cdot 2+7=7$ m/s.
			 \itemch Gia tốc của vật tại thời điểm $t=2$ là $a(2)=v'(2)=s''(2)=6\cdot 2-6=6$ m/s$^2$.
			\itemch Vận tốc của chuyển động bằng $16$ m/s  tại thời điểm $t$ nghĩa là
			$$v(t)=s'(t)=16\Leftrightarrow  3t^2-6t+7=16\Leftrightarrow\hoac{&t=3&&\text{(thỏa mãn)}\\
			 &t=-1&&\text{(loại).}}$$
			Gia tốc của vật tại thời điểm $t=3$ là $a(3)=v'(3)=s''(3)=6\cdot 3-6=12$ m/s$^2$. 
			\itemch Vận tốc của vật  là $v(t)=3t^2-6t+7=3(t-1)^2+4\geq 4,\ \forall t$.\\
			Vậy tại thời điểm $t=1$ thì vận tốc của chuyển động đạt giá trị nhỏ nhất bằng  $4$ m/s.
	\end{itemchoice}}
\end{ex}

% nguồn: Ok_dungsai_lop11_Lan2
\begin{ex}%[1D7V3-3]%[ma:c8511a76]
	Một vật chuyển động theo qui luật là $s=-\dfrac{1}{2}t^3+3t^2+20$ với $t$  (giây) là khoảng thời gian tính từ khi vật bắt đầu chuyển động và $s$ (mét) là quãng đường vật di chuyển được trong khoảng thời gian đó.
	 \choiceTF
		 {Vận tốc của chuyển động bằng $0$ khi $t=0$ s hoặc $t=3$ s}
		{\True Vận tốc của  chuyển động tại thời điểm $t=2$ s là $6$ m/s}
		{\True  Gia tốc của chuyển động bằng $0$ m/s$^2$ khi $t=2$ s}
		{Quãng đường vật đi được bắt đầu từ lúc vật chuyển động tới thời điểm vật đạt được vận tốc  lớn nhất bằng $30$ m}
	\loigiai{Ta có $v(t)=s'(t)=-\dfrac{3}{2}t^2+6t$ và $a(t)=v'(t)=-3t+6$. 
		\begin{itemchoice}
			\itemch Ta  có $v(t)=0\Leftrightarrow -\dfrac{3}{2}t^2+6t-9\Leftrightarrow\hoac{&t=0\\
			 &t=4.}$
			\itemch Khi $t=2$ s thì vận tốc của chuyển động  là $v=-\dfrac{3}{2}\cdot 2^2+6=6$ m/s.
			\itemch  Gia tốc của chuyển động bằng $0$ m/s$^2$ khi
			 $$a(t)=0\Leftrightarrow -3t+6=0\Leftrightarrow t=2\ s.$$
			 \itemch Ta có $v(t)=-\dfrac{3}{2}\left(t^2-4t\right)=-\dfrac{3}{2}\left(t^2-4t+4\right)+6=-\dfrac{3}{2}(t-2)^2+6\leq 6$.\\
			Vậy vận tốc đạt giá trị lớn  nhất bằng $6$ m/s khi $t=2$ s.\\
			Quãng đường vật đi được bắt đầu từ lúc vật chuyển động tới thời điểm vật đạt được vận tốc lớn nhất là $s=-\dfrac{1}{2}\cdot 2^3+3\cdot 2^2+20=28$ m. 
		\end{itemchoice}
	 }
\end{ex}


%%%%%%%%%%%%%%%% TRẢ LỜI NGẮN %%%%%%%%%%%%%%%%

% nguồn: Ok_TraLoiNgan_Lop11_Lan1
\begin{ex}%[1D7H3-3]%[ma:253bca54]
	Một chất điểm chuyển động có phương trình $S=S(t)=t^3+4t^2-2$. Trong đó $t>0$, tính bằng giây $(\mathrm{s})$ và $S$ tính bằng $(\mathrm{m})$. Tính gia tốc của chuyển động  tại thời điểm $t=2\mathrm{\,s}$.
	 \shortans{20}
	\loigiai{ 
		Ta có vận tốc tức thời của chuyển động tại thời  điểm $t$ là $v(t)=S'(t)=3t^2+8t$.\\
		Gia tốc tức  thời của chuyển động tại thời điểm $t$ là $a(t)=v'(t)=6t+8$.\\
		Suy ra gia tốc của chuyển động tại thời điểm $t=2\mathrm{\,s}$  là $a(2)=6\cdot 2+8=20\mathrm{\,m}/\mathrm{s}^2$.
	} 
\end{ex}

% nguồn: Ok_TraLoiNgan_Lop11_Lan1
\begin{ex}%[1D7H3-3]%[ma:643f7857]
	Một chuyển động thẳng xác định bởi phương trình $s(t)=-2t^3+75t+3$, trong đó $s$ tính bằng mét và $t$ là thời gian tính bằng giây. Tính vận tốc và gia tốc  của chuyển động tại thời điểm $t=3$ (s).
	\shortans{-36} 
	 \loigiai{
		Ta có $v(t)=s'(t)=-6t^2+75;\,a(t)=s''(t)=-12t$.\\ 
		Khi đó vận tốc và gia tốc của  chuyển động tại thời điểm $t=3$ là $v(3)=21$ (m/s) và $a(3)=-36$ (m/s$^2$).
	 }
\end{ex}

% nguồn: Ok_TraLoiNgan_Lop11_Lan2
\begin{ex}%[1D7H3-3]%[ma:e8ab3d8d]
	Một vật chuyển động có phương trình $s=t^4-3t^3-3t^2+2t+1$ (m), $t$ là thời gian tính bằng giây. Tính gia tốc của vật tại  thời điểm $t=3$ s.
	\shortans{48} 
	 \loigiai{
		Ta  có $v(t)=s'(t)=4t^3-9t^2-6t+2$ và
		 $a(t)=s''(t)=12t^2-18t-6$.\\
		Gia tốc của vật tại thời điểm $t=3$ s là $a(3)=12\cdot 3^2-18\cdot  3-6=48$ m/s$^2$.
	 }
\end{ex}

% nguồn: Ok_TraLoiNgan_Lop11_Lan2
\begin{ex}%[1D7H3-3]%[ma:9012992d]
	Một vật chuyển động với phương  trình $s(t)=4t^2+t^3$, trong đó $t > 0$, $t$ tính bằng giây, $s(t)$ tính bằng $m$. Tìm gia tốc của vật tại thời điểm vận tốc của vật bằng $ 11 $.
	\shortans{74} 
	\loigiai{ 
		Ta có  $ v(t)=s'(t)=8t+3t^2 $ và $ a(t)=v'(t)=8+6t $.\\
		Khi  đó $ a(11)=8+6\cdot 11= 74$.
	 }
\end{ex}

% nguồn: Ok_TraLoiNgan_Lop11_Lan2
\begin{ex}%[1D7H3-3]%[ma:efd05337]
	Một chất điểm chuyển động có phương trình $s=-t^3+t^2+t+4$ ($t$ là thời gian tính bằng giây). Gia tốc của chuyển động tại thời điểm vận tốc đạt giá trị lớn nhất  là bao nhiêu?
	\shortans{0} 
	\loigiai{ 
		Vận tốc của chất  điểm là $v(t)=s'(t)=-3t^2+2t+1$.\\
		Ta có $v(t)$ là một hàm số bậc hai  có hệ số $a=-3<0$ nên $v(t)$ đạt giá trị lớn nhất khi $t=\dfrac{-b}{2a}=\dfrac{1}{3}$.\\
		 Gia tốc của chất điểm là $s''(t)=-6t+2$.\\
		Tại thời điểm vận tốc đạt giá trị lớn nhất thì gia tốc bằng $s''\left(\dfrac{1}{3}\right)=0$. 
	 }
\end{ex}

% nguồn: Ok_TraLoiNgan_Lop11_Lan2
\begin{ex}%[1D7H3-3]%[ma:3ec9f517]
	Một vật chuyển động với phương trình $s(t)=4t^2+t^3$ , trong  đó $t > 0$, $t$ tính bằng giây, $s(t)$ tính bằng $m$. Tìm gia tốc của vật tại thời điểm vận tốc của vật bằng $ 11 $.
	\shortans{14} 
	 \loigiai{
		Ta có $v(t)=s'(t)=3t^2+8t$ và  $a(t)=v'(t)=6t+8$.\\
		Theo đề bài, ta có $v=11\Rightarrow 3t^2+8t=11\Rightarrow t=1$ $(t >  0)$.\\
	Vậy gia tốc của vật tại thời điểm vận tốc của vật bằng $  11 $ là $a(1)=14m/{s^2}$.}
\end{ex}

% nguồn: Ok_TraLoiNgan_Lop11_Lan2
\begin{ex}%[1D7H3-3]%[ma:a8c5a5e2]
	Một chất điểm chuyển động theo phương trình $s(t)=3\sin 2t+2\cos 2t$, trong đó $t$ là thời gian tính bằng giây và $s$ là quãng đường chuyển động được của chất điểm trong $t$ giây tính bằng mét. Tính gia tốc của chất điểm đó  khi $t=\dfrac{\pi}{4}$.
	 \shortans{-12}
	 \loigiai{
		Ta có \begin{itemize}
		\item $v(t)=s'(t)=3(\sin 2t)'+2(\cos 2t)'=6\cos  2t-4\sin 2t$,
		\item $a(t)=s''(t)=6(\cos 2t)'-4(\sin 2t)'=-12\sin 2t-8\cos  2t$.
	\end{itemize}
	 Gia tốc của chất điểm tại thời điểm $t=\dfrac{\pi}{4}$ là
	$$a\left(\dfrac{\pi}{4}\right)=s''\left(\dfrac{\pi}{4}\right)=-12\left[\sin\left(2\cdot\dfrac{\pi}{4}\right)\right]-8\left[\cos\left(2\cdot\dfrac{\pi}{4}\right)\right]=-12.$$ 
} 
\end{ex}

% nguồn: Ok_TraLoiNgan_Lop11_Lan2
\begin{ex}%[1D7H3-3]%[ma:d77aa7e9]
	Một vật dao động điều hòa có phương trình $x=2\sin\pi t$ ($ x $ tính bằng cm, $t$ tính bằng giây). Tính thời điểm đầu tiên vật có  gia tốc lớn nhất.
	\shortans{1,5} 
	 \loigiai{
		Gia tốc của vật bằng đạo hàm cấp hai của phương  trình chuyển động.\\
		Ta có $x'=2\pi\cos\pi t\Rightarrow a(t)=x''(t)=-2\pi ^2\sin\pi t$.\\ 
		Vì  $-1\le\sin\pi t\le 1\Leftrightarrow-2\pi ^2\le-2\pi ^2\sin\pi t\le 2\pi ^2\Leftrightarrow-2\pi ^2\le a(t)\le 2\pi ^2$.\\
		 Gia tốc lớn nhất khi $\sin\pi t=-1\Leftrightarrow\pi t=-\dfrac{\pi}{2}+k2\pi\Leftrightarrow t=-\dfrac{1}{2}+2k$.\\
		Với $k=0\Rightarrow t=-\dfrac{1}{2}(l);k=1\Rightarrow t=\dfrac{3}{2}$.\\ 
	Vậy tại $t=\dfrac{3}{2}$ giây là  thời điểm đầu tiên vật có gia tốc lớn nhất.}
\end{ex}

% nguồn: Ok_TraLoiNgan_Lop11_Lan2
\begin{ex}%[1D7V3-3]%[ma:e38dc1e8]
	\immini{Một chuyển động có vận tốc được biểu diễn theo đồ thị là một đường parabol (xem hình bên). Tính gia tốc của chuyển động tại thời điểm $t=1$ s.}{\begin{tikzpicture}[scale=1,font=\footnotesize,line join=round,line cap=round,>=stealth,x=0.6cm,y=0.6cm,declare function = {f (\x) =-4*(\x)^2+11*(\x)+3;}]
		 \def \xmin{-1.5};
		\def \xmax{4.8}; 
		\def  \ymin{-1.5};
		\def  \ymax{11.5};
		\draw[->] (\xmin,0.) -- (\xmax,0.) node[anchor=north] {$t$}; 
		\draw[->]  (0.,\ymin) -- (0.,\ymax) node[anchor=west] {$v (t)$};
		\clip (\xmin+0.1,\ymin+0.1) rectangle  (\xmax-0.1,\ymax-0.1);
		\draw[smooth,samples=100,domain=0:3]  plot (\x,{f (\x)});
		\draw[dashed] (2,0) node[below]{$2$} |-  (0,9) node[left]{$9$};
		\foreach \x/\y/\l/\g in {0/0/O/225,3/0/3/-90,0/3/3/180,2/9//0}{\fill[black] (\x,\y) circle  (1pt)+(\g:0.5) node{$\l$};}
	\end{tikzpicture}
	 }
	\shortans{3} 
	\loigiai{ 
		 Đồ thị của vận tốc là một Parabol có phương trình $v(t)=at^2+bt+c$ ($a\neq 0$).\\
		Trên hình vẽ đồ thị qua các điểm $(0;3)$, $(2;9)$, $(3;0)$ nên ta  có hệ phương trình
		 $$\heva{&c=3\\&4a+2b+c=9\\&9a+3b+c=0}\Leftrightarrow\heva{&a=-4\\&b=11\\&c=3.}$$
		Do đó phương trình của vận  tốc là $v(t)=-4t^2+11t+3$.\\
		 Suy ra gia tốc của chuyển động là $a(t)=v'(t)=-8t+11$.\\
		Vậy gia tốc của chuyển động tại thời  điểm $t=1$ (s) là $a(1)=v'(1)=3$.
	 }
\end{ex}

% nguồn: Ok_TraLoiNgan_Lop11_Lan2
\begin{ex}%[1D7V3-3]%[ma:ba647522]
	Một vật chuyển động có vận tốc (m/s) được cho bởi hàm số bậc hai có đồ thị như hình vẽ sau 
	\begin{center}
		\begin{tikzpicture}[scale=1,font=\footnotesize,line join=round,line cap=round,>=stealth,declare function = {f (\x) =2*(\x)^2-4*(\x)-1;}]
				 \def \xmin{-1.5};
				 \def \xmax{2.8};
				 \def \ymin{-3.8};
				 \def \ymax{1.5};
				\draw[->] (\xmin,0.) --  (\xmax,0.) node[anchor=north] {$x$};
				\draw[->] (0.,\ymin)  -- (0.,\ymax) node[anchor=west] {$y$};
				\clip (\xmin+0.1,\ymin+0.1) rectangle  (\xmax-0.1,\ymax-0.1);
				\draw[smooth,samples=100,domain=\xmin:\xmax]  plot (\x,{f (\x)});
				\draw[dashed] (1,0) node[above]{$1$} |- (0,-3)  node[left]{$-3$};
				\foreach \x/\y/\l/\g  in {0/0/O/45,1/0//0,0/-3//0,0/-1/-1/180,1/-3//0}{\fill[black] (\x,\y) circle (1pt)+(\g:0.3) node{$\l$};}
			\end{tikzpicture}
	\end{center}
	Biết gia tốc của vật tại thời điểm vận tốc bị triệt tiêu là $a\sqrt{b}$ m/s$^2$, trong đó $a,b\in\mathbb{N}^*$ và  $b\leq 10$. Tính $2a+3b$.
	\shortans{22} 
	\loigiai{ 
		Gọi hàm số bậc  hai cần xác định là $y=ax^2+bx+c$, ($a\neq 0$).\\
		Tọa độ  đỉnh $I(1;-3)\Rightarrow\dfrac{-b}{2a}=1\Leftrightarrow2a+b=0$.\\
		Điểm $I(1;-3)$ thuộc đồ thị hàm số  nên $a+b+c=-3$.\\
		 Đồ thị cắt $Oy$ tại $(0;-1)\Rightarrow c=-1$.\\
		Do đó  ta có hệ phương trình $\heva{&2a+b=0\\&a+b+c=-3\\&c=-1}\Leftrightarrow\heva{&a=2\\&b=-4\\&c=-1.}$\\
		Như thế,  vận tốc của vật là $v(x)=2x^2-4x-1$.\\
		Ta  có
		$v(x)=0\Leftrightarrow 2x^2-4x-1=0\Leftrightarrow\hoac{&x=\dfrac{2+\sqrt{6}}{2}&&\\  &x=\dfrac{2-\sqrt{6}}{2}&& \text{(loại).}}$\\
		Gia tốc của vật  là $a(x)=4x-4$.\\
		Gia tốc của vật  tại thời điểm vận tốc triệt tiêu là $a\left(\dfrac{2+\sqrt{6}}{2}\right)=4\cdot\dfrac{2+\sqrt{6}}{2}-4=2\sqrt{6}$.\\
		Vậy $a=2$, $b=6$  và $2a+3b=4+18=22$.
	} 
\end{ex}


%%%%%%%%%%%%%%%% TỰ LUẬN %%%%%%%%%%%%%%%%

% nguồn: Ok_TuanLuan_Toan11_Lan2
\begin{ex}%[1D7H3-3]%[ma:09873255]
	Một chuyển động xác định bởi phương trình $s(t)=t^3-3t^2-9t+2$. Trong đó $t$ được tính bằng giây, $S$ được tính bằng mét.  Tính gia tốc của chuyển động tại thời điểm $t=3$ s?
	\loigiai{ 
		Vận tốc của chuyển động tại thời điểm $t$ có phương trình  là $v(t)=s'(t)=3t^2-6t-9$.\\
		Gia tốc của chuyển động tại  thời điểm $t$ có phương trình là $a(t)=v'(t)=6t-6$.\\
		Tại thời điểm $t=3$ s ta có  $a(3)=6\cdot 3-6=12$ m/s$^2$.
	 }
\end{ex}

% nguồn: Ok_TuanLuan_Toan11_Lan2
\begin{ex}%[1D7H3-3]%[ma:1a985217]
	Một chất điểm chuyển động có phương trình $S=2t^4+6t^2-3t+1$ với $t$ tính bằng giây (s) và $S$ tính bằng mét (m). Hỏi gia tốc của chuyển động tại thời điểm $t=3$ (s) bằng bao nhiêu? 
	\loigiai{ 
		Ta có  vận tốc tức thời của chuyển động được tính theo công thức
		$v(t)=s'(t)=8t^3+12t-3$.\\ 
		Khi đó gia tốc tức  thời của chuyển động được tính theo công thức
		$a(t)=v'(t)=24t^2+12$.\\ 
		Vậy gia tốc của chuyển động tại thời điểm $t=3$ (s) là $a(3)=24\cdot 3^2+12=228$ m/s$^2$. 
	} 
\end{ex}

% nguồn: Ok_TuanLuan_Toan11_Lan2
\begin{ex}%[1D7H3-3]%[ma:f8b05299]
	Cho chuyển động thẳng xác định bởi phương trình $S=-t^3+3t^2+9t$, trong đó $t$ tính bằng giây và $S$ tính bằng mét. Tính vận tốc của chuyển động tại thời điểm gia tốc triệt  tiêu.
	\loigiai{ 
		Vận tốc của chuyển động chính là đạo hàm cấp một  của quãng đường $v=S'=-3t^2+6t+9$.\\
		Gia tốc của chuyển động chính là đạo hàm cấp  hai của quãng đường $a=S''=-6t+6$.\\
		Gia tốc triệt tiêu khi  $S''=0\Leftrightarrow t=1$.\\
		Khi đó vận tốc của chuyển động là $S'(1)=12$  m/s.
	 }
\end{ex}
